\documentclass[10pt,a4paper]{amsart}
\usepackage[margin=19mm]{geometry}
\usepackage{amsmath,amssymb,mathtools}
\usepackage[scr=rsfs,cal=euler]{mathalfa}
\usepackage{xfrac}
\usepackage[unicode,hidelinks,bookmarks=false]{hyperref}
\newcommand{\ie}{i.e.\ }
\newcommand{\cf}{cf.\ }
\newcommand{\resp}{respectively\ }
\newcommand{\Subsection}{\subsection}
\providecommand{\nobreakdash}{\nobreak\hbox{-}}
\setlength{\emergencystretch}{2em}
\multlinegap0pt
\usepackage{bm}
\usepackage{bbm}
\usepackage{dsfont}
\usepackage{xspace}
\usepackage{cancel}
\usepackage{mathtools}
\usepackage{amsthm}
\makeatletter
\@ifundefined{theorem}{\newtheorem{theorem}{Theorem}}{}
\@ifundefined{proposition}{\newtheorem{proposition}[theorem]{Proposition}}{}
\makeatother
\DeclareMathOperator{\Aut}{Aut}
\newcommand \bff {{\mathbf{f}}}
\newcommand \bfg {{\mathbf{g}}}
\newcommand \calO {{\mathcal{O}}}
\newcommand \kx {\Bbbk^\times}
\newcommand{\SB}{\mathbb{S}}
\title[Poisson fields of two variables]{Poisson fields of two variables}
\author{Ken Goodearl, James J. Zhang}
\date{Source: arXiv:2605.24835v1 (CC BY 4.0)}
\begin{document}
\maketitle

\noindent\emph{Source and licence.} Poisson fields of two variables. Version arXiv:2605.24835v1 (CC BY 4.0), distributed under the Creative Commons Attribution 4.0 International licence, as recorded in the arXiv licence metadata for this exact version and shown on the abstract page https://arxiv.org/abs/2605.24835v1. Full rights notice: licenses/arxiv-2605.24835-CC-BY-4.0.txt.\par\smallskip

\noindent\textit{Editorial scope.} Counted unit: the Theorem whose source label is \texttt{zzthm7.5} in the article (source0.tex lines 3350--3375, source label \texttt{zzthm7.5}), delivered together with the article's own complete original proof of it (source0.tex lines 3377--3467, one \texttt{proof} environment of 91 lines). No mathematics is written, repaired or re-derived: statement and proof are the article's own text line for line. Required context retained uncounted: zzpro7.3 (lines 3306--3334). Every definition, display and label of those lines is retained unchanged. Omitted from this package and not redistributed: 1--3305, 3335--3349, 3376--3376, 3468--4016. Nothing inside a retained excerpt is omitted or altered: every retained body is delivered exactly as the article's lines are written, so no omission receipt of any kind accompanies this package. Citations: the retained excerpts carry no citation command at all, so no bibliography entry of the article is transcribed and no reference is redistributed. Reference adaptation: none was needed and none was made; every reference inside the delivered unit resolves inside it, so no unresolved reference or citation is produced. Printed slips kept literal: no printed word, symbol, hypothesis or number of the counted statement or of its proof was altered. Layout and class: the article's own class is \texttt{amsart} with packages including latexsym, amsxtra, amscd, ifthen, amsfonts, verbatim, amsmath, amsthm, amssymb, url; the wrapper is the lane's pinned \texttt{amsart} editorial wrapper at 10pt on A4, which restates the theorem-like environments the retained text needs and adds the article's own macro definitions that the retained text needs (\texttt{\textbackslash Aut}, \texttt{\textbackslash SB}, \texttt{\textbackslash bff}, \texttt{\textbackslash bfg}, \texttt{\textbackslash calO}, \texttt{\textbackslash kx}), with extra wrapper packages (bm, bbm, dsfont, xspace, cancel, mathtools). The delivered numbering is the unit's own. Assets: no retained span contains a figure environment, a graphics-inclusion command, a PDF-page inclusion, a table or a TikZ picture; no image file is copied into this package, the package declares no asset dependency and no image file is redistributed with it. Licence: version arXiv:2605.24835v1 of this article is distributed under the Creative Commons Attribution 4.0 International licence, as recorded in the arXiv licence metadata for this exact version and shown on the abstract page https://arxiv.org/abs/2605.24835v1. A full rights notice is delivered at licenses/arxiv-2605.24835-CC-BY-4.0.txt. Characters: every retained excerpt is pure ASCII, so the delivered engine, XeLaTeX with its default Latin Modern text font, reports no missing character.
\begin{proposition}
\label{zzpro7.3}
Let $\bff$ and $\bfg$ be two nonzero polynomials in $\Bbbk[x,y]$. Suppose that
\begin{enumerate}
\item[(a)]
$\phi$ is a Poisson morphism from $K\{\bff\}$ to $K\{\bfg\}$, and
\item[(b)]
$\bff$ is of the form 
$c\prod_{i=1}^m (x+\xi_i) \prod_{j=1}^{n} (y+\chi_j)$
where $c \in \kx$ and $(\xi_i)_{i=1}^{m}$, $(\chi_j)_{j=1}^n$ 
are families of distinct scalars in $\Bbbk$, with $m,n\ge3$.
\end{enumerate}
Then the following hold.
\begin{enumerate}
\item[(1)]
$\phi(\bff) \in \Bbbk[x,y] \bfg$ and $\deg \bff\leq \deg \bfg$.  
\item[(2)]
If $\deg \bff=\deg \bfg$, then  $\phi$ preserves the degree.
Consequently, $\phi$ is a $\Bbbk$-algebra isomorphism sending 
$x\mapsto c_{11}x+c_{12}y+c_1$ and $y\mapsto c_{21}x+c_{22}y+c_2$ 
for some $c_1,c_2\in \Bbbk$ and some $\begin{pmatrix} c_{11}&c_{12}\\
c_{21}&c_{22}\end{pmatrix} \in GL_2(\Bbbk)$. Further, $\phi(\bff) = (c_{11}c_{22} - c_{12}c_{21}) \bfg$.
\item[(3)]{\rm{(Dixmier Problem)}}
$K\{\bff\}$ has the Dixmier property.
\item[(4)]
$\Aut_{Poi}(K\{\bff\})$ is isomorphic to a subgroup of $\{\phi\in \Aut(\Bbbk[x,y])
\mid \phi(\bff) \in \kx \bff\}.$
\end{enumerate}
\end{proposition}

\begin{theorem}[Automorphism Problem]  \label{zzthm7.5} 
Let  $\bff = c\prod_{i=1}^m (x+\xi_i) \prod_{j=1}^{n} (y+\chi_j)$
where $c \in \kx$ and $(\xi_i)_{i=1}^{m}$, $(\chi_j)_{j=1}^n$ 
are families of distinct scalars in $\Bbbk$ and $m,n \ge 3$.  Define
\begin{itemize}
\item 
$G := \{ \phi \in \Aut_{Poi}(K\{\bff\}) \mid \phi(x) \in \Bbbk[x],\; \phi(y) \in \Bbbk[y] \}$;
\item 
$m^* := m-1$ if there is some $s \in [1,m]$ such that $m \xi_s = \sum_{i=1}^m \xi_i$, and $m^* := m$ otherwise;
\item 
$n^* := n-1$ if there is some $s \in [1,n]$ such that $n \chi_s = \sum_{j=1}^n \chi_j$, and $n^* := n$ otherwise;
\item 
$H := \{ (\gamma,\delta) \in (\kx)^2 \mid \gamma^{m^*} = \delta^{n^*} = \gamma^{m-1} \delta^{n-1} = 1 \}$.
\end{itemize}
Then the following hold.
\begin{enumerate}
\item[(1)] 
$G$ is a normal subgroup of $\Aut_{Poi}(K\{\bff\})$ with index at most $2$.
\item[(2)]
$G$ is isomorphic to a subgroup of $H$.
\item[(3)]
$\Aut_{Poi}(K\{\bff\}) = G$ except possibly when $m=n$ is even and $m^* = n^* = m$, in which case $G$ is cyclic of order at most $m$.
\item[(4)]
$|{\Aut_{Poi}}(K\{\bff\})| \le (m-1)(n-1)$.
\end{enumerate}
\end{theorem}

\begin{proof}
(1) We first claim that any $\phi \in \Aut_{Poi}(K\{\bff\})$ is either in $G$ or satisfies $\phi(x) \in \Bbbk[y]$, $\phi(y) \in \Bbbk[x]$.  By Proposition \ref{zzpro7.3}(2), we have $\phi(x) = c_{11}x+c_{12}y+c_1$ and $\phi(y) = c_{21}x+c_{22}y+c_2$ 
for some $c_1,c_2\in \Bbbk$ and some $\bigl( \begin{smallmatrix} c_{11}&c_{12}\\
c_{21}&c_{22}\end{smallmatrix} \bigr) \in GL_2(\Bbbk)$, and $\phi(\bff) \in \kx \bff$.  Thus,
\begin{equation}  \label{E7.5.1}  \tag{E7.5.1}
\phi(\bff) = c \prod_{i=1}^m (c_{11}x+c_{12}y+c_1 +\xi_i) \prod_{j=1}^n (c_{21}x+c_{22}y+c_2 + \chi_j).
\end{equation}
Since $\phi(\bff) \in \kx \bff$, the factor $c_{11}x+c_{12}y+c_1 +\xi_1$ must be a scalar multiple of some $x+\xi_i$ or of some $y+\chi_j$, and likewise for $c_{21}x+c_{22}y+c_2 + \chi_1$.  Note that if $\bigl( d_{ij} \bigr) = \bigl( c_{ij} \bigr)^{-1}$, then
\begin{align*}
d_{11}(c_{11}x+c_{12}y+c_1 +\xi_1) + d_{12}(c_{21}x+c_{22}y+c_2 + \chi_1) &= x + \text{a scalar}  \\
d_{21}(c_{11}x+c_{12}y+c_1 +\xi_1) + d_{22}(c_{21}x+c_{22}y+c_2 + \chi_1) &= y + \text{a scalar},
\end{align*}
so $c_{11}x+c_{12}y+c_1 +\xi_1$ and $c_{21}x+c_{22}y+c_2 + \chi_1$ cannot both be in $\Bbbk[x]$ or both in $\Bbbk[y]$.  Thus, either $c_{12} = c_{21} = 0$, in which case $\phi \in G$, or else $c_{11} = c_{22} = 0$, whence $\phi(x) \in \Bbbk[y]$, $\phi(y) \in \Bbbk[x]$.  This verifies the claim.

Observe from \eqref{E7.5.1} that we can only have $c_{11} = c_{22} = 0$ when $m=n$.  Thus, $\Aut_{Poi}(K\{\bff\}) = G$ when $m \ne n$.

It is clear that $G$ is a subgroup of $\Aut_{Poi}(K\{\bff\})$, and it follows from the claim that $G$ is normal.  Suppose there exists $\psi \in \Aut_{Poi}(K\{\bff\}) \setminus G$.  For any $\phi \in \Aut_{Poi}(K\{\bff\}) \setminus G$, we see from the claim that $\phi\psi^{-1} \in G$, whence $\phi G = \psi G$.  Consequently, 
$$
|{\Aut_{Poi}}(K\{\bff\}) / G | \le 2.
$$

(2) For any $\phi \in G$, we have, by definition,
$$
\phi(x) = c_{11}(\phi) x + c_1(\phi) \qquad \text{and} \qquad \phi(y) = c_{22}(\phi) y + c_2(\phi)
$$
for some $c_{11}(\phi), c_{22}(\phi) \in \kx$ and $c_1(\phi), c_2(\phi) \in \Bbbk$.  The rule $\phi \mapsto (c_{11}(\phi), c_{22}(\phi))$ defines a homomorphism $\Phi : G \rightarrow (\kx)^2$.  We need to show that $\Phi$ is injective and $\Phi(G) \subseteq H$.

Fix $\phi \in G$, and abbreviate $c_{ii} := c_{ii}(\phi)$, $c_i := c_i(\phi)$.  In view of Proposition \ref{zzpro7.3}(2), \eqref{E7.5.1} reduces to
\begin{equation}  \label{E7.5.2}  \tag{E7.5.2}
c_{11} c_{22} \bff = \phi(\bff) = c \prod_{i=1}^m (c_{11}x+c_1 +\xi_i) \prod_{j=1}^n (c_{22}y+c_2 + \chi_j).
\end{equation}
Comparing leading coefficients in this equation, we obtain $c_{11} c_{22} c = c c_{11}^m c_{22}^n$, and so $c_{11}^{m-1} c_{22}^{n-1} = 1$.  We will be done on showing that
\begin{enumerate}
\item[(a)]
$c_{11}^{m^*} = 1$, and $c_{11} = 1 \implies c_1 = 0$;
\item[(b)]
$c_{22}^{n^*} = 1$, and $c_{22} = 1 \implies c_2 = 0$.
\end{enumerate}
Since (a) and (b) are symmetric, we just prove (a).

Comparing factors from $\Bbbk[x]$ in \eqref{E7.5.2}, we see that there must be a permutation $\sigma \in \SB_m$ such that $c_{11}x + c_1 + \xi_i = c_{11}(x + \xi_{\sigma(i)})$ for all $i$, that is,
\begin{equation}  \label{E7.5.3}  \tag{E7.5.3}
c_{11} \xi_{\sigma(i)} = c_1 + \xi_i \qquad \forall\; i \in [1,m].
\end{equation}
Summing this equation over $i \in [1,m]$ yields 
\begin{equation}  \label{E7.5.4}  \tag{E7.5.4}
(c_{11}-1) \sum_{i=1}^m \xi_i = m c_1 \,.
\end{equation}
In particular, the implication $c_{11}  = 1 \implies c_1 = 0$ follows.

Since $1^{m^*} = 1$, we assume $c_{11} \ne 1$ for the remainder of the proof of (a).

\textbf{Claim 1:} $s \in [1,m]$ is a fixed point of $\sigma$ if and only if $\xi_s = c_1/(c_{11}-1)$, if and only if $m \xi_s = \sum_{i=1}^m \xi_i$.

The first forward implication follows from \eqref{E7.5.3}, and the second equivalence is immediate from \eqref{E7.5.4}.  If $\xi_s = c_1/(c_{11}-1)$, then $c_{11} \xi_s = c_1 + \xi_s = c_{11} \xi_{\sigma(s)}$, and so $\sigma(s) = s$ because the $\xi_i$ are distinct.

\textbf{Claim 2:} $\xi_{\sigma^t(i)} = (c_{11}^{-1} + \cdots + c_{11}^{-t}) c_1 + c_{11}^{-t} \xi_i$ for all $i \in [1,m]$ and $t \ge 0$.

This is an obvious induction on $t$.

\textbf{Claim 3:} Let $\calO$ be a $\sigma$-orbit of cardinality $t>1$.  Then $c_{11}$ is a primitive $t$-th root of unity, i.e., $t$ is the order of $c_{11}$ in the group $\kx$.

If $i \in \calO$, then $\sigma^t(i) = i$, and by Claim 2 we see that $(c_{11}^t - 1) \xi_i = \frac{c_{11}^t-1}{c_{11}-1}\, c_1$.  Since there are at least $2$ indices in $\calO$, we must have $c_{11}^t = 1$.  Thus, $c_{11}$ has order $d<\infty$ in $\kx$, and $d \mid t$.  Now $c_{11}^d \xi_i = (c_{11}^{d-1} + \cdots + 1) c_1 + \xi_i$ for any $i \in \calO$.  Claim 2 shows that $\xi_{\sigma^d(i)} = \xi_i$ and hence $\sigma^d(i) = i$.  This forces $t=d$, proving Claim 3.

We can now complete the proof of (a) and thus of (2).  If $\sigma$ has a fixed point $s$, then by Claim 1, $s$ is the only fixed point of $\sigma$, and $m^* = m-1$.  By Claim 3, all the non-singleton orbits of $\sigma$ have the same cardinality $t$, and $c_{11}^t = 1$.  Since $[1.m] \setminus \{s\}$ is a disjoint union of these latter orbits, $t \mid m-1$, and thus $c_{11}^{m^*} = 1$.

If $\sigma$ has no fixed points, then $m^* = m$ and $[1,m]$ is a disjoint union of $\sigma$-orbits of cardinality $t$.  In this case, $c_{11}^t = 1$ and $t \mid m$, so again $c_{11}^{m^*} = 1$.

(3) That $\Aut_{Poi}(K\{\bff\}) = G$ when $m \ne n$ was already noted in the proof of part (1).  Suppose there exists $\theta \in \Aut_{Poi}(K\{\bff\}) \setminus G$, so that $m=n$.  Now $\theta(x) = c_{12}y+c_1$ and $\theta(y) = c_{21}x+c_2$ for some $c_{12},c_{21} \in \kx$ and $c_1,c_2 \in \Bbbk$.  Observe that $\theta^2 \in G$ and
$$
\theta^2(x) = c_{12}c_{21} x + c_{12}c_2+c_1 \,, \qquad\qquad \theta^2(y) = c_{12}c_{21} y + c_{21}c_1+c_2 \,.
$$
Then $(c_{12}c_{21}, c_{12}c_{21}) \in H$, whence
$$
(c_{12}c_{21})^{m^*} = (c_{12}c_{21})^{n^*} = (c_{12}c_{21})^{2(m-1)} = 1.
$$

In the present case, \eqref{E7.5.1} reads
\begin{align*}
\theta(\bff) &= c \prod_{i=1}^m (c_{12}y+c_1+\xi_i) \prod_{j=1}^m (c_{21}x+c_2+\chi_j)  \\
&= (c_{12}c_{21})^m c \prod_{j=1}^m (x+c_{21}^{-1}(c_2+\chi_j)) \prod_{i=1}^m (y+c_{12}^{-1}(c_1+\xi_i)) \,,
\end{align*}
from which we conclude that $\theta(\bff) = (c_{12}c_{21})^m \bff$.  On the other hand, Proposition \ref{zzpro7.3}(2) says that $\theta(\bff) = - c_{12}c_{21} \bff$.  Consequently, $(c_{12}c_{21})^{m-1} = -1$.  In particular, $m^*,n^* \ne m-1$, and so $m^*=n^* = m$.  Thus $c_{12}c_{21} = -1$ and $m$ is even.

We now get $H = \{ (\gamma,\gamma^{-1}) \in (\kx)^2 \mid \gamma^m = 1 \}$, which is cyclic of order at most $m$, and then the same holds for $G$ by part (2).

(4)  We first show that $|H| \le (m-1)(n-1)$, which is clear in case $m^* = m-1$ and $n^* = n-1$.
If $m^* = m$, then all $(\gamma,\delta) \in H$ satisfy $\gamma = \delta^{n-1}$, whence $|H| \le n \le (m-1)(n-1)$.  Similarly, $|H| \le m \le (m-1)(n-1)$ in case $n^* = n$.  Therefore, we have $|{\Aut_{Poi}}(K\{\bff\})| \le (m-1)(n-1)$ whenever $\Aut_{Poi}(K\{\bff\}) = G$.

Finally, suppose that $\Aut_{Poi}(K\{\bff\}) \ne G$, whence $G$ has index $2$ and $|{\Aut_{Poi}}(K\{\bff\})| = 2 |G|$.  By part (3), $m=n$ is even and $|G| \le m$.  Since then $m \ge 4$, we conclude that $|{\Aut_{Poi}}(K\{\bff\})| \le 2m < (m-1)^2$ in this case.
\end{proof}

\end{document}
